\begin{theorem}[激活分支的三阶接触刚性与唯一小根性]\label{thm:xi-output-potential-activated-branch-rigidity}
\leanverified{paper\_xi\_output\_potential\_activated\_branch\_rigidity}
沿用命题 \ref{prop:real-input-40-activated-branch-linear} 的记号，令
\[
G(\lambda,u)=\lambda^8F(\lambda^{-1},u).
\]
则在 \(u=1\) 的某邻域内，激活分支 \(\lambda_{\mathrm{act}}(u)\) 满足：
\begin{enumerate}
  \item \textbf{（三阶接触刚性）}
  \[
  \boxed{\;
  \lambda_{\mathrm{act}}(u)-(u-1)=3(u-1)^3+O\!\bigl((u-1)^4\bigr).
  \;}
  \]
  换言之，\(\lambda_{\mathrm{act}}\) 对线性分支 \(u\mapsto u-1\) 恰有二阶切触，而首个非线性偏离出现在三次项。
  \item \textbf{（唯一小根性）} 存在解析函数 \(\widetilde G(\lambda,u)\)，使
  \[
  G(\lambda,u)=\bigl(\lambda-\lambda_{\mathrm{act}}(u)\bigr)\widetilde G(\lambda,u),
  \qquad
  \widetilde G(0,1)\neq 0.
  \]
  因而在 \(u=1\) 的某个复邻域内，\(\lambda_{\mathrm{act}}(u)\) 是唯一趋于 \(0\) 的根分支，而其余七条根分支都与 \(0\) 保持正距离。
\end{enumerate}
\end{theorem}
\begin{proof}
由注 \ref{rem:real-input-40-activated-branch-series}，
\[
\lambda_{\mathrm{act}}(u)
:=(u-1)+3(u-1)^3-(u-1)^4+18(u-1)^5-22(u-1)^6
+O\!\bigl((u-1)^7\bigr).
\]
因此二次项系数严格为零，而三次项系数等于 \(3\neq 0\)，从而立刻得到第一条。

再由命题 \ref{prop:real-input-40-activated-branch-linear}，\(\lambda=0\) 是 \(u=1\) 时的简单根，且
\[
\lambda_{\mathrm{act}}(1)=0,
\qquad
\partial_\lambda G(0,1)=1\neq 0.
\]
由隐函数定理，存在唯一解析分支 \(\lambda_{\mathrm{act}}(u)\) 满足
\[
G(\lambda_{\mathrm{act}}(u),u)=0
\]
并经过 \((\lambda,u)=(0,1)\)。

对解析函数 \(G(\lambda,u)\) 在 \(\lambda=\lambda_{\mathrm{act}}(u)\) 处应用 Weierstrass 除法，可写成
\[
G(\lambda,u)=\bigl(\lambda-\lambda_{\mathrm{act}}(u)\bigr)\widetilde G(\lambda,u),
\]
其中 \(\widetilde G\) 在 \((0,1)\) 邻域解析。把 \(\lambda=\lambda_{\mathrm{act}}(u)\) 处求导代入上式可得
\[
\partial_\lambda G(\lambda_{\mathrm{act}}(u),u)=\widetilde G(\lambda_{\mathrm{act}}(u),u).
\]
因此在 \((0,1)\) 处，
\[
\widetilde G(0,1)=\partial_\lambda G(0,1)=1\neq 0.
\]
由连续性可知，存在 \(\delta,r>0\)，使得当 \(|u-1|<\delta\) 且 \(|\lambda|<r\) 时，
\[
\widetilde G(\lambda,u)\neq 0.
\]
于是 \(G(\lambda,u)=0\) 在该邻域内只能由
\(
\lambda=\lambda_{\mathrm{act}}(u)
\)
给出。这说明 \(\lambda_{\mathrm{act}}(u)\) 是唯一小根，而其余七条根分支全部远离 \(0\)。证毕。
\end{proof}

\paragraph{单位根扭曲的固定模态展开与四阶修正增益}

\begin{theorem}[固定单位根模态误差指数的四阶展开]\label{thm:xi-output-potential-dirichlet-fixed-j-fourth-order}
\leanverified{paper\_xi\_output\_potential\_dirichlet\_fixed\_j\_fourth\_order}
沿用命题 \ref{prop:real-input-40-output-cumulants-closed} 与命题 \ref{prop:real-input-40-output-dirichlet-nonrh} 的记号。令
\[
\omega_m:=e^{2\pi i/m},
\qquad
\rho_{m,j}:=\rho\!\bigl(M(\omega_m^j)\bigr),
\qquad
\theta_{m,j}:=\frac{\log \rho_{m,j}}{\log \lambda},
\qquad
\lambda=\lambda(1)=\varphi^2,
\]
并记
\[
\kappa_2:=P''(0)>0,
\qquad
\kappa_4:=P^{(4)}(0)>0.
\]
则对每个固定整数 \(j\ge 1\)，当 \(m\to\infty\) 且 \(m>j\) 时有
\[
\log\frac{\rho_{m,j}}{\lambda}
:=
-\frac{\kappa_2}{2}\Bigl(\frac{2\pi j}{m}\Bigr)^2
+\frac{\kappa_4}{24}\Bigl(\frac{2\pi j}{m}\Bigr)^4
+O_j(m^{-6}),
\]
从而
\[
\boxed{\;
\theta_{m,j}
:=
1-\frac{2\pi^2\kappa_2}{\log\lambda}\frac{j^2}{m^2}
+\frac{2\pi^4\kappa_4}{3\log\lambda}\frac{j^4}{m^4}
+O_j(m^{-6}).\;}
\]
特别地，
\[
\boxed{\;
\theta_{m,1}
:=
1-\frac{2\pi^2\kappa_2}{\log\lambda}\,m^{-2}
+\frac{2\pi^4\kappa_4}{3\log\lambda}\,m^{-4}
+O(m^{-6}).\;}
\]
\end{theorem}
\begin{proof}
由命题 \ref{prop:real-input-40-output-cumulants-closed}，压力
\[
P(\theta)=\log \lambda(e^\theta)
\]
在 \(\theta=0\) 邻域解析，且各阶导数均为实数。另一方面，命题 \ref{prop:real-input-40-output-dirichlet-nonrh} 给出的单位根扭曲满足
\[
\log\frac{\rho_{m,j}}{\lambda}
:=
\Re\!\left(P\!\left(\frac{2\pi i j}{m}\right)-P(0)\right).
\]
把 \(P(\theta)-P(0)\) 在 \(\theta=0\) 处展开为 Taylor 级数，并取
\(
\theta=2\pi i j/m
\)，得到
\[
P(\theta)-P(0)
:=
P'(0)\theta+\frac{\kappa_2}{2}\theta^2+\frac{P^{(3)}(0)}{6}\theta^3
+\frac{\kappa_4}{24}\theta^4+O_j(m^{-5}).
\]
由于奇次项在纯虚方向上贡献纯虚部，而首个未写出的偶次实部出现在六阶，故
\[
\Re\!\left(P\!\left(\frac{2\pi i j}{m}\right)-P(0)\right)
:=
-\frac{\kappa_2}{2}\Bigl(\frac{2\pi j}{m}\Bigr)^2
+\frac{\kappa_4}{24}\Bigl(\frac{2\pi j}{m}\Bigr)^4
+O_j(m^{-6}).
\]
这就给出第一式。再除以 \(\log \lambda\) 即得 \(\theta_{m,j}\) 的展开。\(j=1\) 情形只是其特例。证毕。
\end{proof}

\begin{corollary}[有界单位根窗口内的最坏模态最终由 \texorpdfstring{$j=1$}{j=1} 实现]\label{cor:xi-output-potential-dirichlet-bounded-window-j1}
\leanverified{paper\_xi\_output\_potential\_dirichlet\_bounded\_window\_j1}
表 \ref{tab:real_input_40_output_potential_dirichlet_twists} 在已审计模数上记录了
\[
j_{\mathrm{worst}}=1.
\]
更进一步，对任意固定整数 \(J\ge 2\)，存在 \(m_0(J)\) 使得对所有 \(m\ge m_0(J)\) 与一切 \(2\le j\le J\)，均有
\[
\theta_{m,1}>\theta_{m,j}.
\]
因此，在任何有界 Dirichlet 扭曲窗口内，\(j=1\) 最终严格支配全部其余固定模态。
\end{corollary}
\begin{proof}
令
\[
A:=\frac{2\pi^2\kappa_2}{\log\lambda}>0,
\qquad
B:=\frac{2\pi^4\kappa_4}{3\log\lambda}>0.
\]
由定理 \ref{thm:xi-output-potential-dirichlet-fixed-j-fourth-order}，对每个固定 \(J\ge 2\)，可把 \(2\le j\le J\) 上的误差统一写成 \(O_J(m^{-6})\)，从而
\[
\theta_{m,1}-\theta_{m,j}
:=
A\frac{j^2-1}{m^2}
-B\frac{j^4-1}{m^4}
+O_J(m^{-6}).
\]
右端主项为正且阶为 \(m^{-2}\)；其余项至多为 \(m^{-4}\)。因此对每个固定 \(j\in\{2,\dots,J\}\)，当 \(m\) 足够大时差值严格为正。对有限集合 \(\{2,\dots,J\}\) 取所需阈值的最大者，即得结论。证毕。
\end{proof}

\noindent 上述推论只使用了 \(j\) 有界的局部 Taylor 窗口。再结合单位圆上谱半径函数在 \(t=0\) 邻域的严格下降与远离 \(1\) 时的统一谱隙，还可把表 \ref{tab:real_input_40_output_potential_dirichlet_twists} 中的经验模式提升为充分大模数下的全局严格结论。
